% Generated draft: requires mathematical and source review.
\documentclass[11pt]{article}
\usepackage[margin=25mm]{geometry}
\usepackage{fontspec}
\setmainfont{Latin Modern Roman}
\usepackage{amsmath,amssymb,mathtools,unicode-math}
\setmathfont{Latin Modern Math}
\usepackage{enumitem,needspace,xurl,hyperref}
\setlength{\parindent}{0pt}
\setlength{\parskip}{5pt}
\setlength{\emergencystretch}{3em}
\newcommand{\N}{\mathbb{N}}
\newcommand{\Z}{\mathbb{Z}}
\newcommand{\Q}{\mathbb{Q}}
\newcommand{\R}{\mathbb{R}}
\newcommand{\C}{\mathbb{C}}
\newcommand{\F}{\mathbb{F}}
\newcommand{\A}{\mathbb{A}}
\newcommand{\PP}{\mathbb{P}}
\newcommand{\E}{\mathbb{E}}
\newcommand{\eps}{\varepsilon}
\newcommand{\dd}{\,\mathrm{d}}
\begin{document}
\section*{UnsolvedMath ID: 30004532}
\textbf{Source problem code:} catalogue entry
\subsection{Definitions and mathematical statement}
Let \(k\) be a field. In the usual commutative graded interpretation, a standard graded Artinian \(k\)-algebra is a finite-dimensional commutative algebra
\[
A=\bigoplus_{i\geq 0}A_i,\qquad A_0=k,
\]
generated as a \(k\)-algebra by \(A_1\). Put \(\mathfrak m=\bigoplus_{i>0}A_i\). Such an algebra is Gorenstein if its socle
\[
\operatorname{Soc}(A)=\{a\in A: \mathfrak m a=0\}
\]
has dimension one over \(k\). Its codimension is \(\dim_k A_1\). If \(s=\max\{i:A_i\neq 0\}\), its \(h\)-vector is
\[
h(A)=(h_0,\ldots,h_s),\qquad h_i=\dim_k A_i.
\]
This vector is unimodal if there exists \(t\in\{0,\ldots,s\}\) such that
\[
h_0\leq\cdots\leq h_t\geq h_{t+1}\geq\cdots\geq h_s,
\]
with empty chains at the endpoints interpreted vacuously; it is non-unimodal otherwise.

For a specified field \(k\), the existence question is: does there exist a standard graded Artinian Gorenstein \(k\)-algebra \(A\) satisfying \(\dim_k A_1=4\) whose \(h\)-vector is non-unimodal?

The supplied source does not specify the base field, its characteristic, or whether the field is fixed or may be chosen as part of the example. Consequently, this formulation specifies the question relative to \(k\), while the intended quantifier and restrictions on \(k\) remain unresolved.
\subsection{Short English statement}
Can a standard graded Artinian Gorenstein algebra of codimension four have a Hilbert function that fails to increase weakly to a peak and then decrease weakly? The supplied statement leaves the base field and its characteristic unspecified.
\subsection{Sources}
\begin{itemize}
\item Local UnsolvedMath registry, supplied record with permanent ID \(30004532\), source problem code \texttt{catalogue entry}, in the Oberwolfach Reports Open Problems collection. No individual proposer is identified.
\item Supplied report citation: \emph{Lefschetz Properties in Algebra, Geometry and Combinatorics} (2021), Oberwolfach Reports, pp.\ 1539--1600, DOI \texttt{10.4171/owr/2020/31}. \url{https://doi.org/10.4171/owr/2020/31}. The registry separately records the proposed year as 2020.
\item Seo--Srinivasan, supplied literature reference: \url{https://arxiv.org/abs/1011.2732}.
\item Migliore--Nagel--Zanello, supplied literature reference: \url{https://intlpress.com/site/pub/files/_fulltext/journals/mrl/2008/0015/0002/MRL-2008-0015-0002-a011.pdf}.
\item Literature assessment dated exactly \texttt{2026-08-22}: the supplied curation labels the problem open and reports that existence of a non-unimodal Gorenstein \(h\)-vector in codimension four remains unresolved, with partial results constraining possible vectors but neither producing a counterexample nor proving universal unimodality. This is supplied dated curation derived from the OWR material, not an independently verified literature assessment; its status and cited claims require independent verification.
\end{itemize}
\end{document}
